% =============================================================================
% RESENHA CRÍTICA — MODELO ABNT
% =============================================================================
% Características da resenha crítica:
%   - Síntese da obra (1/3 do texto)
%   - Análise crítica fundamentada (2/3 do texto)
%   - Identificação completa da obra no cabeçalho
%   - Destinada a periódico científico ou trabalho acadêmico
% =============================================================================
\documentclass[
  12pt,
  a4paper,
  brazilian
]{article}

% --- Pacotes ---
\usepackage[utf8]{inputenc}
\usepackage[T1]{fontenc}
\usepackage[brazilian,portuges]{babel}
\usepackage[alf]{abntex2cite}
\usepackage{abntex2abrev}
\usepackage{setspace}
\usepackage{geometry}
\usepackage{hyperref}
\usepackage{microtype}

% --- Margens ABNT ---
\geometry{
  a4paper,
  left=3cm,
  top=3cm,
  right=2cm,
  bottom=2cm
}

\onehalfspacing

\hypersetup{
  pdfauthor={Nome do Autor},
  pdftitle={Resenha: [Título da Obra]},
  colorlinks=true,
  linkcolor=black,
  citecolor=black,
  urlcolor=black
}

\title{Resenha de \textit{Título da Obra Resenhada}}
\author{Nome do Resenhista \\ \small{Instituição} \\ \small{e-mail}}
\date{2026}

\begin{document}

\maketitle

\section*{Identificação da Obra}

\begin{flushleft}
\textbf{Obra:} Título completo da obra \\
\textbf{Autor(es):} Nome do autor \\
\textbf{Editora:} Nome da editora \\
\textbf{Cidade:} Cidade de publicação \\
\textbf{Ano:} 2026 \\
\textbf{Número de páginas:} XXX p. \\
\textbf{ISBN:} 978-XX-XXXX-XXX-X
\end{flushleft}

\section*{Síntese da Obra}

A presente resenha tem por objetivo apresentar e analisar criticamente a
obra \textit{Título da Obra}, de Nome do Autor. A obra está estruturada em
[X] capítulos, distribuídos em [X] partes, nos quais o autor aborda...

[Parágrafo 1: contextualização do tema e do autor]

[Parágrafo 2: estrutura da obra — capítulos, seções, organização interna]

[Parágrafo 3: resumo do conteúdo de cada capítulo/parte]

\section*{Análise Crítica}

\subsection*{Contribuições da Obra}

Discussão sobre os méritos da obra: originalidade, consistência teórica,
rigor metodológico, relevância, atualidade.

\subsection*{Limites e Lacunas}

Apontamento crítico de aspectos que poderiam ser aprofundados,
controvérsias não resolvidas, limitações metodológicas ou teóricas.

\subsection*{Diálogo com a Literatura}

Comparação com outras obras que tratam do mesmo tema, situando a
contribuição específica da obra resenhada.

\section*{Considerações Finais}

Avaliação geral da obra (recomendada ou não, para qual público, em que
contextos), explicitando o valor da contribuição e as limitações
identificadas.

\section*{Referências}

\bibliographystyle{abntex2-alf}
\bibliography{resenha_modelo}

\end{document}
